% ===============================================================
%  Chapter 6 -- A History of the Protocols               (Part I)
% ===============================================================
\chapter{A History of the Protocols}
\label{ch:history}

\begin{figure}[htbp]
\centering
\adjustbox{max width=\linewidth}{%
\begin{forest}
  chapter map,
  [Chapter 6
    [First schemes [Directory 1978] [Write-once 1983]]
    [Snooping in the 1980s [Illinois {\&} Berkeley] [MOESI vocabulary 1986]]
    [Beyond the bus [{DASH, tokens}] [{AMD MOESI, Intel MESIF}]]
  ]
\end{forest}}
\caption[]{Structure of Chapter~6: the protocols of
Chapter~5 in the order they appeared.}
\label{fig:ch6-map}
\end{figure}

The protocols of Chapter~5 grew out of work on bus-based multiprocessors in the
1980s, long before multicore chips. \Cref{fig:timeline} shows the milestones that
matter for this report, and Part~III returns to several of them.

\begin{figure}[htbp]
\centering
\begin{bigfig}
\begin{tikzpicture}[font=\sffamily\scriptsize,
    ev/.style={draw, thick, rounded corners=2pt, align=center, inner sep=2pt,
               text width=2.1cm}]
  \draw[very thick, slateC, -{Stealth}] (-0.6,0) -- (13.9,0);
  \foreach \i/\y/\lev/\txt/\col in {
      0/1978/1.15/{Directory scheme\\Censier {\&} Feautrier}/amberC,
      1/1983/-1.15/{Write-once\\Goodman}/skyC,
      2/1984/2.55/{Illinois protocol (MESI)\\Papamarcos {\&} Patel}/mesiBlue,
      3/1986/-2.55/{MOESI states\\Sweazey {\&} Smith}/moesiOrange,
      4/1986/1.15/{Six protocols compared\\Archibald {\&} Baer}/plumC,
      5/1990/-1.15/{DASH directory\\Lenoski et al.}/amberC,
      6/2001/2.55/{First dual-core chip\\IBM POWER4}/tealC,
      7/2003/-2.55/{Token coherence\\Martin et al.}/brickC,
      8/2005/1.15/{Intel and AMD\\go dual-core}/tealC,
      9/2009/-1.15/{MESIF report\\Goodman {\&} Hum}/roseC,
      10/2012/2.55/{``Coherence is\\here to stay''}/leafC} {
    \pgfmathsetmacro{\x}{\i*1.32}
    \pgfmathsetmacro{\e}{\lev>0 ? \lev-0.42 : \lev+0.42}
    \draw[\col, thick] (\x,0) -- (\x,\e);
    \fill[\col] (\x,0) circle (2.4pt);
    \node[ev, draw=\col, fill=\col!9] at (\x,\lev) {\textbf{\y}\\\txt};
  }
\end{tikzpicture}
\end{bigfig}
\caption[Timeline of coherence protocols]{Milestones from the first directory
scheme to current multicore chips, in date order but not to scale. Sources are
cited in the text.}
\label{fig:timeline}
\end{figure}

% ---------------------------------------------------------------
\section{The First Coherence Schemes}
\label{sec:hist-before}
\label{sec:hist-first}
\sectionmap{[First schemes [Several processors on one bus] [1978: a central directory]
  [1983: write-once on a bus]]}
\index{snooping protocol}\index{directory protocol}%

Multiprocessors go back to the 1960s and 1970s, and by the late 1990s four or
eight processors commonly shared one motherboard and one
bus~\cite[\S12.1]{sarangi2025basic}. As soon as processors with private caches
shared a bus, coherence became an engineering problem. The first general scheme
was a directory. Censier and Feautrier proposed in 1978 that memory keep, for each
block, a record of the caches holding it, so that invalidations went only where
they were needed~\cite{censier1978directory}. The first widely known snooping
protocol was Goodman's write-once~\cite{goodman1983writeonce}, with the states
Invalid, Valid, Reserved and Dirty~\cite{archibald1986evaluation}. The first write
to a block went through to memory and invalidated the other copies, and later
writes stayed in the cache. Every later protocol keeps the main idea of
write-once, which is to use the bus both for moving data and for telling other
caches what has happened to it.

% ---------------------------------------------------------------
\section{Snooping Protocols of the 1980s and the MOESI Vocabulary}
\label{sec:hist-1980s}
\label{sec:hist-moesi}
\sectionmap{[1980s snooping [Illinois: the E state] [Berkeley: a shared-dirty owner]
  [Firefly and Dragon: update] [1986: five states describe them all]]}
\index{MESI}\index{MOESI}\index{MSI}\index{MESIF}\index{Owned state}\index{Exclusive state}\index{snooping protocol}\index{write-update}\index{Berkeley protocol}\index{Dragon protocol}\index{Firefly protocol}\index{Illinois protocol}\index{Futurebus}\index{owner}%

In 1986 Archibald and Baer described and simulated six snooping protocols side by
side~\cite{archibald1986evaluation}. \Cref{tab:hist-protocols} translates their
state names into the terms of Chapter~5. Two of them matter most here. The
Illinois protocol of Papamarcos and Patel added a Valid-Exclusive state so that a
later write needed no bus transaction~\cite{papamarcos1984illinois}, which is
MESI. The Berkeley protocol had a Shared-Dirty state in which one cache kept
modified data, other caches held copies, and the Shared-Dirty cache supplied the
data to readers~\cite{archibald1986evaluation}. That is the \stO\ state before it
had the name. Firefly and Dragon were write-update protocols.

\begin{table}[htbp]
\centering
\caption{The six protocols evaluated by Archibald and
Baer~\cite{archibald1986evaluation}, with each original state name mapped to
its nearest MOESI equivalent.}
\label{tab:hist-protocols}
\small
\begin{tabular}{@{}l l l l@{}}
\toprule
Protocol & Original states & In MOESI terms & Strategy \\
\midrule
Write-once & Invalid, Valid, Reserved, Dirty & I, S, E, M & invalidate \\
Synapse    & Invalid, Valid, Dirty & I, S, M & invalidate \\
Illinois   & Invalid, Valid-Exclusive, Shared, Dirty & I, E, S, M & invalidate \\
Berkeley   & Invalid, Valid, Shared-Dirty, Dirty & I, S, \textbf{O}, M & invalidate \\
Firefly    & Valid-Exclusive, Shared, Dirty & E, S, M & update \\
Dragon     & Valid-Exclusive, Shared-Clean, Shared-Dirty, Dirty & E, S, \textbf{O}, M & update \\
\bottomrule
\end{tabular}
\end{table}

By the mid-1980s these protocols used different names for overlapping ideas.
While working on the IEEE Futurebus standard, Sweazey and Smith picked out five
states (Modified, Owned, Exclusive, Shared and Invalid) that could describe a
whole class of protocols, and showed how protocols of that class could share one
bus~\cite{sweazey1986moesi}. MOESI began as a shared vocabulary for existing
protocols more than as one new design. MESI and the ownership idea each add one
state to MSI, MOESI combines the two, and MESIF (Part~III) extends MESI in another
direction.

% ---------------------------------------------------------------
\section{Beyond the Bus and onto the Chip}
\label{sec:hist-scaling}
\label{sec:hist-multicore}
\sectionmap{[Recent history [1990: DASH directory] [2003: token coherence]
  [AMD: MOESI; Intel: MESIF]]}
\index{MOESI}\index{MSI}\index{MESIF}\index{directory protocol}\index{token coherence}\index{DASH}%

A bus carries one message at a time, which limits the number of processors it can
serve. The Stanford DASH machine of 1990 kept MSI-family states but replaced the
bus with a directory and a network of point-to-point links~\cite{lenoski1990dash}.
In 2003 Martin, Hill and Wood proposed token coherence~\cite{martin2003token}.
Both come back in Part~III. When the cores moved onto one chip in the
2000s~\cite[\S12.1]{sarangi2025basic}, AMD documented a MOESI
protocol~\cite{amd64vol2} and Intel's QuickPath processors adopted
MESIF~\cite{goodman2009mesif}. In 2012 Martin, Hill and Sorin argued that hardware
coherence will stay practical as core counts grow~\cite{martin2012here}. AMD and
Intel have kept different extensions of the same protocol for years, and Part~II
asks what each extension costs.
